% Generated draft: requires mathematical and source review.
\documentclass[11pt]{article}
\usepackage[margin=25mm]{geometry}
\usepackage{fontspec}
\setmainfont{Latin Modern Roman}
\usepackage{amsmath,amssymb,mathtools,unicode-math}
\setmathfont{Latin Modern Math}
\usepackage{enumitem,needspace,xurl,hyperref}
\setlength{\parindent}{0pt}
\setlength{\parskip}{5pt}
\setlength{\emergencystretch}{3em}
\newcommand{\N}{\mathbb{N}}
\newcommand{\Z}{\mathbb{Z}}
\newcommand{\Q}{\mathbb{Q}}
\newcommand{\R}{\mathbb{R}}
\newcommand{\C}{\mathbb{C}}
\newcommand{\F}{\mathbb{F}}
\newcommand{\A}{\mathbb{A}}
\newcommand{\PP}{\mathbb{P}}
\newcommand{\E}{\mathbb{E}}
\newcommand{\eps}{\varepsilon}
\newcommand{\dd}{\,\mathrm{d}}
\begin{document}
\section*{UnsolvedMath ID: 30003626}
\textbf{Source problem code:} catalogue entry
\subsection{Definitions and mathematical statement}
Use the conventional finite-gap interpretation
\[
E=\bigcup_{j=1}^{m}[\alpha_j,\beta_j],\qquad m\geq1,\qquad \alpha_1<\beta_1<\alpha_2<\cdots<\alpha_m<\beta_m.
\]
Write \(\operatorname{dist}(z,S)=\inf_{s\in S}|z-s|\). Interpret \(\partial E\) as the boundary in \(\R\), namely the band endpoints \(\{\alpha_j,\beta_j:1\leq j\leq m\}\). These conventions require confirmation against the cited report; the supplied record does not explicitly define them.

Under the usual interpretation, \(\mathcal T_E\) is the isospectral torus of bounded, two-sided, self-adjoint Jacobi operators with spectrum \(E\) that are reflectionless on \(E\). Here reflectionlessness means that, for every \(n\in\Z\), the diagonal Green function satisfies
\[
\operatorname{Re}\left(\lim_{\eta\downarrow0}\langle e_n,(J_0-x-i\eta)^{-1}e_n\rangle\right)=0
\quad\text{for Lebesgue-almost every }x\in E,
\]
where \((e_n)_{n\in\Z}\) is the standard basis of \(\ell^2(\Z)\). The record supplies the notation \(\mathcal T_E\), but not its definition.

Let \(J_0\in\mathcal T_E\), with
\[
(J_0u)_n=a^{(0)}_{n-1}u_{n-1}+b^{(0)}_nu_n+a^{(0)}_nu_{n+1},
\qquad a^{(0)}_n>0,\quad b^{(0)}_n\in\R.
\]
Let \(J=J_0+\delta J\) be a bounded complex Jacobi operator on \(\ell^2(\Z)\), with compact \(\delta J\), written using the indexing convention
\[
(Ju)_n=c_{n-1}u_{n-1}+b_nu_n+a_nu_{n+1},
\qquad a_n,b_n,c_n\in\C.
\]
Set \(\delta a_n=a_n-a^{(0)}_n\), \(\delta b_n=b_n-b^{(0)}_n\), and \(\delta c_n=c_n-a^{(0)}_n\). No self-adjointness condition is imposed on \(J\). Write \(\sigma_d(J)\) for its isolated eigenvalues of finite algebraic multiplicity, and count eigenvalues with algebraic multiplicity in the sums below. This multiplicity convention is standard but is not stated explicitly in the supplied record.

For the exponent \(p\) occurring in the cited estimate, put
\[
S_p(\delta J)=\sum_{n\in\Z}\bigl(|\delta a_n|^p+|\delta b_n|^p+|\delta c_n|^p\bigr).
\]
\textbf{Missing hypothesis:} the supplied record does not specify the admissible range of \(p\). Thus the following questions are conditional on the exponent hypotheses of the cited bound; they do not assert validity for every real or every positive \(p\). The substantive estimates concern perturbations for which \(S_p(\delta J)<\infty\); compactness alone does not supply this summability assumption.

The supplied reconstruction of bound (7) is
\[
\sum_{z\in\sigma_d(J)}
\frac{\operatorname{dist}(z,E)^{p+\varepsilon}(1+|z|)^{1/2-2\varepsilon}}
{\operatorname{dist}(z,\partial E)^{1/2}}
\leq L_{\varepsilon,p,E}S_p(\delta J).
\tag{7}
\]
Its constant is understood to depend only on the indicated parameters, uniformly over \(J_0\in\mathcal T_E\) and the admissible compact perturbations.

\textbf{Endpoint question.} For each admissible \(p\) and finite-gap set \(E\), does there exist a finite constant \(L_{0,p,E}\) such that every such \(J_0\) and perturbation satisfy
\[
\sum_{z\in\sigma_d(J)}
\frac{\operatorname{dist}(z,E)^p(1+|z|)^{1/2}}
{\operatorname{dist}(z,\partial E)^{1/2}}
\leq L_{0,p,E}S_p(\delta J)?
\]

\textbf{Denominator question.} Under the intended parameter hypotheses of (7), can its denominator be replaced by \(\operatorname{dist}(z,E)^{1/2}\)? Equivalently, does there exist a finite constant \(\widetilde L_{\varepsilon,p,E}\), with the same uniformity, such that
\[
\sum_{z\in\sigma_d(J)}
\operatorname{dist}(z,E)^{p+\varepsilon-1/2}(1+|z|)^{1/2-2\varepsilon}
\leq\widetilde L_{\varepsilon,p,E}S_p(\delta J)?
\]
The record does not determine whether this second question concerns positive \(\varepsilon\), the endpoint \(\varepsilon=0\), or both. These alternatives must remain distinct pending source clarification. The supplied dated curation describes (7) as proved for positive \(\varepsilon\) and the two proposed improvements as unresolved; these are supplied assessments, not independently verified mathematical or literature conclusions.
\subsection{Short English statement}
For compact complex perturbations of a reflectionless finite-gap Jacobi operator, ask whether the stated eigenvalue estimate holds at \(\varepsilon=0\), and whether distance to the band endpoints in its denominator can be replaced by distance to the whole spectrum. The admissible range of \(p\) and the intended \(\varepsilon\)-range of the second question are missing. This is a statement-only draft requiring human mathematical review.
\subsection{Sources}
\begin{itemize}
\item Local UnsolvedMath registry: supplied record for permanent ID \(30003626\), source code \texttt{catalogue entry}, title ``Sharp Eigenvalue Bounds for Nonselfadjoint Jacobi Perturbations,'' in the Oberwolfach Reports Open Problems collection. Project registry locations are \texttt{data/unsolvedmath/problems.json} and the supplementary \texttt{catalog/problems.json}; no independent file inspection is claimed.
\item Supplied source citation: \emph{Mini-Workshop: Reflectionless Operators: The Deift and Simon Conjectures} (2018), Oberwolfach Reports, pp.\ 2943--2985, DOI \texttt{10.4171/owr/2017/49}. \url{https://doi.org/10.4171/owr/2017/49}. The record assigns proposed year 2017 and names no individual proposer. Its original statement refers to inequality (7).
\item Supplied literature source: Christiansen--Zinchenko. \url{https://doi.org/10.1007/s11005-017-0961-z}. No fuller bibliographic description is supplied.
\item Supplied literature source: Oberwolfach Report 49/2017. \url{https://doi.org/10.4171/owr/2017/49}.
\item Literature assessment dated exactly \texttt{2026-08-21}: the supplied curation labels the problem open, reports the finite-gap complex perturbation estimate as proved for positive \(\varepsilon\), and describes the endpoint and denominator replacement as unresolved. This is supplied dated curation derived from the OWR record, not independently verified; time-sensitive status and cited claims require independent checking.
\item The registry labels the statement \texttt{corrected\_verified} and says that the official report was checked to reconstruct the full setup and inequality (7). This is a supplied verification claim, not verification performed for this draft.
\end{itemize}
\end{document}
